% GENERATED FILE. Edit canonical lessons, then run npm run generate:books.
% canonical-source-hash: fnv1a64:4e5b4861e1c42fc3
% canonical-lessons: TA-C446-kanakku, TA-C446-iruppu, TA-C446-vatti, TA-C446-kacolai, TA-C446-rokkam

\chapter{At the Bank}
\label{ch:ta-at-the-bank}
\hlchaptermodality{446}

\begin{chapteropening}
\textbf{By the end of this chapter:} \emph{I can talk about a bank account and its balance, interest, cheques and cash.}

Complete the last lesson of chapter 446: I can talk about a bank account and its balance, interest, cheques and cash.
\end{chapteropening}

\section[\texorpdfstring{\ta{கணக்கு} (kaṇakku)}{கணக்கு (kaṇakku)}]{\ta{கணக்கு} (kaṇakku) — an account; a calculation}

\label{lesson:TA-C446-kanakku}



\begin{quote}
Before the new one: say the Tamil for difficulty, then the Tamil for to shake.
\end{quote}

\subsection*{You'll want to know: \ta{கணக்கு}}
\textbf{\ta{கணக்கு}} — \emph{kaṇakku} — \textquotedblleft{}an account; a calculation\textquotedblright{}.

\begin{quote}
\ta{நான்} \ta{வங்கியில்} \ta{ஒரு} \ta{கணக்கு} \ta{திறந்தேன்}.
\end{quote}

\begin{quote}
\emph{nāṉ vaṅkiyil oru kaṇakku tiṟantēṉ.}
\end{quote}

\begin{quote}
I opened an account at the bank.
\end{quote}

Also a sum in arithmetic: \ta{கணக்குப்} \ta{போடு}, to work out a sum.

\begin{cousinweb}
\textbf{\ta{கணக்கு}} (\emph{kaṇakku}) is usually traced to the Sanskrit family of \emph{gaṇana}, \textquotedblleft{}counting\textquotedblright{} (CDIAL 3992).
\end{cousinweb}

\begin{grammarlens}[title={What you've built}]
One more word for getting around.
\end{grammarlens}

\subsection*{Your turn}
\begin{itemize}
\raggedright
  \item \emph{Say it:} \emph{kaṇakku}
  \item \emph{Say it:} \emph{kaṇakku}, once more
  \item \emph{Say it:} \emph{nāṉ vaṅkiyil oru kaṇakku tiṟantēṉ.} — I opened an account at the bank.
  \item \emph{Recall:} say the Tamil for difficulty, then the Tamil for to shake, then say \emph{kaṇakku} again
\end{itemize}

\begin{tcolorbox}[breakable,colback=teal!4,colframe=teal!35!black,title={Before you move on}]
What does \ta{கணக்கு} mean? (\textquotedblleft{}An account; a calculation\textquotedblright{}.) Say it once more.
\end{tcolorbox}
\section[\texorpdfstring{\ta{இருப்பு} (iruppu)}{இருப்பு (iruppu)}]{\ta{இருப்பு} (iruppu) — a balance; stock in hand}

\label{lesson:TA-C446-iruppu}



\begin{quote}
Before the new one: say the Tamil for to shake, then the Tamil for an account.
\end{quote}

\subsection*{You'll want to know: \ta{இருப்பு}}
\textbf{\ta{இருப்பு}} — \emph{iruppu} — \textquotedblleft{}a balance; stock in hand\textquotedblright{}.

\begin{quote}
\ta{என்} \ta{கணக்கில்} \ta{இருப்பு} \ta{குறைவாக} \ta{இருக்கிறது}.
\end{quote}

\begin{quote}
\emph{eṉ kaṇakkil iruppu kuṟaivāka irukkiṟatu.}
\end{quote}

\begin{quote}
The balance in my account is low.
\end{quote}

From \ta{இரு}, to be, + -\ta{ப்பு}: what is there.

\begin{cousinweb}
\textbf{\ta{இருப்பு}} (\emph{iruppu}) is native Dravidian (DEDR 480), the noun of \ta{இரு}, \textquotedblleft{}to be\textquotedblright{}: what remains, a balance in hand.
\end{cousinweb}

\begin{grammarlens}[title={What you've built}]
One more word for getting around.
\end{grammarlens}

\subsection*{Your turn}
\begin{itemize}
\raggedright
  \item \emph{Say it:} \emph{iruppu}
  \item \emph{Say it:} \emph{iruppu}, once more
  \item \emph{Say it:} \emph{eṉ kaṇakkil iruppu kuṟaivāka irukkiṟatu.} — The balance in my account is low.
  \item \emph{Recall:} say the Tamil for to shake, then the Tamil for an account, then say \emph{iruppu} again
\end{itemize}

\begin{tcolorbox}[breakable,colback=teal!4,colframe=teal!35!black,title={Before you move on}]
What does \ta{இருப்பு} mean? (\textquotedblleft{}A balance; stock in hand\textquotedblright{}.) Say it once more.
\end{tcolorbox}
\section[\texorpdfstring{\ta{வட்டி} (vaṭṭi)}{வட்டி (vaṭṭi)}]{\ta{வட்டி} (vaṭṭi) — interest on money}

\label{lesson:TA-C446-vatti}



\begin{quote}
Before the new one: say the Tamil for an account, then the Tamil for a balance.
\end{quote}

\subsection*{You'll want to know: \ta{வட்டி}}
\textbf{\ta{வட்டி}} — \emph{vaṭṭi} — \textquotedblleft{}interest on money\textquotedblright{}.

\begin{quote}
\ta{இந்தக்} \ta{கடனுக்கு} \ta{வட்டி} \ta{அதிகம்}.
\end{quote}

\begin{quote}
\emph{intak kaṭaṉukku vaṭṭi atikam.}
\end{quote}

\begin{quote}
The interest on this loan is high.
\end{quote}

\ta{வட்டிக்குக்} \ta{கொடு} is to lend at interest.

\begin{cousinweb}
\textbf{\ta{வட்டி}} (\emph{vaṭṭi}) is usually traced to Prakrit \emph{vaḍḍhi}, \textquotedblleft{}increase\textquotedblright{}, from Sanskrit \emph{vṛddhi} (CDIAL 12076).
\end{cousinweb}

\begin{grammarlens}[title={What you've built}]
One more word for getting around.
\end{grammarlens}

\subsection*{Your turn}
\begin{itemize}
\raggedright
  \item \emph{Say it:} \emph{vaṭṭi}
  \item \emph{Say it:} \emph{vaṭṭi}, once more
  \item \emph{Say it:} \emph{intak kaṭaṉukku vaṭṭi atikam.} — The interest on this loan is high.
  \item \emph{Recall:} say the Tamil for an account, then the Tamil for a balance, then say \emph{vaṭṭi} again
\end{itemize}

\begin{tcolorbox}[breakable,colback=teal!4,colframe=teal!35!black,title={Before you move on}]
What does \ta{வட்டி} mean? (\textquotedblleft{}Interest on money\textquotedblright{}.) Say it once more.
\end{tcolorbox}
\section[\texorpdfstring{\ta{காசோலை} (kācōlai)}{காசோலை (kācōlai)}]{\ta{காசோலை} (kācōlai) — a cheque}

\label{lesson:TA-C446-kacolai}



\begin{quote}
Before the new one: say the Tamil for a balance, then the Tamil for interest on money.
\end{quote}

\subsection*{You'll want to know: \ta{காசோலை}}
\textbf{\ta{காசோலை}} — \emph{kācōlai} — \textquotedblleft{}a cheque\textquotedblright{}.

\begin{quote}
\ta{அவர்} \ta{காசோலை} \ta{மூலம்} \ta{பணம்} \ta{செலுத்தினார்}.
\end{quote}

\begin{quote}
\emph{avar kācōlai mūlam paṇam celuttiṉār.}
\end{quote}

\begin{quote}
He paid by cheque.
\end{quote}

\ta{காசு}, money, + \ta{ஓலை}, a palm-leaf note.

\begin{cousinweb}
\textbf{\ta{காசோலை}} (\emph{kācōlai}) is a modern coinage built from \ta{காசு}, \textquotedblleft{}a coin, money\textquotedblright{} (DEDR 1431), and \ta{ஓலை}, \textquotedblleft{}a palm leaf for writing\textquotedblright{} (DEDR 1070).
\end{cousinweb}

\begin{grammarlens}[title={What you've built}]
One more word for getting around.
\end{grammarlens}

\subsection*{Your turn}
\begin{itemize}
\raggedright
  \item \emph{Say it:} \emph{kācōlai}
  \item \emph{Say it:} \emph{kācōlai}, once more
  \item \emph{Say it:} \emph{avar kācōlai mūlam paṇam celuttiṉār.} — He paid by cheque.
  \item \emph{Recall:} say the Tamil for a balance, then the Tamil for interest on money, then say \emph{kācōlai} again
\end{itemize}

\begin{tcolorbox}[breakable,colback=teal!4,colframe=teal!35!black,title={Before you move on}]
What does \ta{காசோலை} mean? (\textquotedblleft{}A cheque\textquotedblright{}.) Say it once more.
\end{tcolorbox}
\section[\texorpdfstring{\ta{ரொக்கம்} (rokkam)}{ரொக்கம் (rokkam)}]{\ta{ரொக்கம்} (rokkam) — cash}

\label{lesson:TA-C446-rokkam}



\begin{quote}
Before the new one: say the Tamil for interest on money, then the Tamil for a cheque.
\end{quote}

\subsection*{You'll want to know: \ta{ரொக்கம்}}
\textbf{\ta{ரொக்கம்}} — \emph{rokkam} — \textquotedblleft{}cash\textquotedblright{}.

\begin{quote}
\ta{இந்தக்} \ta{கடையில்} \ta{ரொக்கமாக} \ta{மட்டும்} \ta{பணம்} \ta{வாங்குகிறார்கள்}.
\end{quote}

\begin{quote}
\emph{intak kaṭaiyil rokkamāka maṭṭum paṇam vāṅkukiṟārkaḷ.}
\end{quote}

\begin{quote}
This shop takes payment only in cash.
\end{quote}

\ta{ரொக்கமாக} is in cash.

\begin{cousinweb}
\textbf{\ta{ரொக்கம்}} (\emph{rokkam}), \textquotedblleft{}cash, ready money\textquotedblright{}, is a word whose origin is not settled.
\end{cousinweb}

\begin{grammarlens}[title={What you've built}]
One more word for getting around.
\end{grammarlens}

\subsection*{Your turn}
\begin{itemize}
\raggedright
  \item \emph{Say it:} \emph{rokkam}
  \item \emph{Say it:} \emph{rokkam}, once more
  \item \emph{Say it:} \emph{intak kaṭaiyil rokkamāka maṭṭum paṇam vāṅkukiṟārkaḷ.} — This shop takes payment only in cash.
  \item \emph{Recall:} say the Tamil for interest on money, then the Tamil for a cheque, then say \emph{rokkam} again
\end{itemize}

\begin{tcolorbox}[breakable,colback=teal!4,colframe=teal!35!black,title={Before you move on}]
What does \ta{ரொக்கம்} mean? (\textquotedblleft{}Cash\textquotedblright{}.) Say it once more.
\end{tcolorbox}
